\subsection*{逻辑的形式化}

从我们所拥有的（极不完整的）文本中得到的总体印象是，公元前五世纪的希腊哲学思想被一种日益自觉的努力所主导，即把在修辞学和当代数学中取得巨大成功的清晰论述程序扩展到人类思想的整个领域——换言之，创造最广义上的逻辑。哲学著作的风格在这一时代发生了突变。在公元前七世纪和六世纪，哲学家们断言和预言（或者至多勾勒出基于同样模糊类比的模糊论证），而从巴门尼德，尤其是芝诺开始，他们进行推理，并寻求提取一般原则作为其辩证法的基础。排中律原则的首次陈述见于巴门尼德的著作，而埃利亚的芝诺的“归谬”证明一直享有盛名。但芝诺活跃于公元前五世纪中叶，尽管文献记载不确定（*），但很可能在这一时期数学家们也在自己的领域使用了这些原则。

正如我们已说过的，本书的范围不包括叙述这一逻辑发展各阶段所出现的无数困难，以及由此产生的、曾使埃利亚学派、智者派以及柏拉图和亚里士多德卷入其中的论战。我们只需指出在这一演变过程中，演说术的勤勉培养以及随之自然产生的语言分析所起的作用——这些发展主要归功于公元前五世纪的智者派。另一方面，数学的影响，即使不总是被明确承认，也同样清晰可见，尤其是在柏拉图和亚里士多德的著作中。可以说，柏拉图几乎为数学所痴迷；虽然他本人在这一领域并非创造者，但他始终了解同时代数学家的发现（其中许多人是他的朋友或学生），并且始终对这一学科保持着最直接的兴趣，以至于提出了新的研究方向。在他的著作中，正是数学不断地充当例证或模型（有时，如同毕达哥拉斯学派那样，助长了他对神秘主义的偏爱）。亚里士多德作为柏拉图的学生，几乎不可能没有接受过学园对学生们所要求的最低限度的数学教育，而且一部汇集其著作中涉及或提及数学的选段集已经编纂完成 {[}2 bis{]}。然而，他似乎从未付出多少努力去跟上他那个时代的数学运动，在这一领域他只引用那些早已成为常识的结果。后来的哲学家们，大多对数学的了解甚至更少。他们中许多人缺乏必要的技术知识，他们对数学的引用所涉及的，是这一学科演变中早已过去的阶段。

就逻辑而言，这一时期的顶峰是亚里士多德的巨著 {[}2{]},其伟大功绩在于首次成功地将此前一直模糊不清或未被明确表述的推理方法加以系统化和法典化（*）。就我们的目的而言，我们需要记住这部著作的总论点，即有可能将每一个正确的论证都归结为对少数不变规则的系统应用，这些规则独立于所考察对象的特定性质（这种独立性通过使用字母来表示概念或命题而清楚地显现出来——这一用法亚里士多德很可能借鉴自数学家）。但亚里士多德几乎将全部注意力集中于特定类型的关系和推理链，即他所谓的“三段论”；本质上，这些是用集合论的语言表达的形式为 A ⊂ B 或 A ∩ B ≠ ∅ 的关系（†），以及通过下述图式将这些关系或其否定连接起来的推理链。

\[
(\mathrm{A} \subset \mathrm{B} \text {   且   } \mathrm{B} \subset \mathrm{C}) \Longrightarrow (\mathrm{A} \subset \mathrm{C}).
\]

亚里士多德对他那个时代的数学十分熟悉，不可能意识不到这类图式不足以解释数学家所使用的全部逻辑运算，更不足以解释逻辑在其他方面的应用（{[}2{]}, 《前分析篇》，I，35； {[}2 bis{]}, 第25-26页）（*）。无论如何，他对“三段论”的各种形式进行了细致的研究（几乎完全致力于阐明因推理中所涉及术语的歧义或晦涩而产生的永恒困难），这使他得以制定出，除其他事项外，对命题取否定的规则（{[}2{]}, 《前分析篇》，I，46）。同样归功于亚里士多德的，是他极其清晰地区分了“全称”命题与“特称”命题的作用：这是迈向量词的第一步（†）。但众所周知，他的著作的影响——常常被狭隘而愚蠢地解释——直至十九世纪仍然相当巨大，这导致哲学家们忽视了对数学的研究，并阻碍了形式逻辑的进步（**）。尽管如此，逻辑在古代仍继续取得进展，在麦加拉学派和斯多亚学派——逍遥学派的对手——中。遗憾的是，我们对这些学说的了解全都是第二手的，往往来自其对手或平庸的评注者。看来这些逻辑学家的贡献在于奠定了现代意义上的“命题演算”的基础；他们不像亚里士多德那样局限于特定形式 A⊂B 的命题，而是陈述了关于完全不确定的命题的规则。此外，他们如此细致地分析了这些规则之间的逻辑关系，以至于能够用与我们第一章第2节和第3节所描述的方法非常相似的方法，将它们全部从五条被归类为“不可证明的”基本规则中推导出来 {[}5{]}。遗憾的是，他们的影响转瞬即逝，他们的成果被遗忘，直到十九世纪才被逻辑学家重新发现。亚里士多德在逻辑领域一直是无可争议的权威，直至十七世纪。特别是，经院哲学家们完全处于他的影响之下，尽管他们对形式逻辑的贡献远非微不足道，但与古代哲学家们的贡献相比，其中并不包含第一流的进展。

然而，亚里士多德或其后继者的著作似乎并未在数学领域引起显著反响。希腊数学家们沿着毕达哥拉斯学派及其公元前四世纪后继者（塞奥多洛、泰阿泰德、欧多克索斯）开辟的道路继续研究，在呈现成果时显然并未关注形式逻辑。这并不令人意外，因为将此时期数学推理的灵活性与精确性同亚里士多德逻辑的极其原始状态相比，便可见一斑。而后来，当逻辑超越这一阶段时，其发展又受到了数学新成就的引导。

随着代数学的发展，人们很难不注意到形式逻辑的规则与代数的规则之间的相似性，两者都适用于未指定的对象（命题或数）。而当十七世纪代数符号在韦达和笛卡尔手中定型时，几乎立即出现了各种用符号表示逻辑运算的尝试。但在莱布尼茨之前，这些尝试（例如埃里戈纳（1644年）用符号书写初等几何证明，或佩尔（1659年）对算术做同样的事）都是肤浅的，并未在数学推理分析方面带来任何进展。

莱布尼茨是一位哲学家，同时也是一位一流的数学家，他能够从自己的数学经验中汲取思想的萌芽，这些思想后来使形式逻辑摆脱了经院哲学的困境（*）。如果说有谁堪称旷世奇才，那么莱布尼茨当之无愧；他是一座取之不尽、用之不竭的原创而富有成果的思想宝库。他对逻辑学尤为感兴趣，因为逻辑学正是他毕生致力于形式化语言与思维这一宏伟计划的核心。早在童年时期，他就已与经院逻辑决裂，并被一个可追溯至雷蒙·卢尔的思想所吸引：即把所有人类概念分解为原始概念的方法，这些原始概念由此构成“人类思想的字母表”，并以近乎机械的方式重新组合它们，从而得出所有真命题。{[}12 b{]}，第VII卷，第185页；参见{[}12 bis{]}, 第II章）。当他还很年轻时，他还构想出另一个更具原创性的想法，即符号记法作为思想“阿里阿德涅之线”的效用（*）。“真正的方法，”他说，“应当为我们提供一条阿里阿德涅之线，也就是说，一种确定的、可感知的物质手段，它将引导理智，就像几何中绘制的线条以及算术中为初学者规定的运算形式一样。没有它，我们的心灵将无法在漫长的道路上行走而不迷失方向。”（{[}12 b{]}，第 VII 卷，第 22 页；参见{[}12 bis{]}，第90页）。他在二十五岁之前对当时的数学知之甚少，最初以“通用语言”的形式提出他的计划。但当他接触到代数时，他将其作为他的“通用特征”的模型。通过这个短语，他指的是一种符号语言，能够无歧义地表达所有人类思想，增强我们的推理能力，通过纯粹机械的注意力努力避免错误，并且其构造方式使得“那些连提出者自己都不理解的幻想，将无法用这些符号写出来”（{[}12 a{]}，第I卷，第187页）。在莱布尼茨著作中无数段落里，他提及这一宏伟计划及其实现将带来的进展（参见{[}12 bis{]}，第IV章和第VI章），人们可以看到他多么清晰地构想了一种形式化语言的概念，即作为符号的纯粹组合，其中只有书写顺序是重要的（†），因此机器将能够产生所有定理（***），所有争议都能通过简单计算得到解决（{[}12 bis{]}，第VII卷，第198-203页）。尽管这些希望可能显得夸张，但不可否认的是，莱布尼茨思想中的这一恒定主题支撑了他相当一部分数学成果，从他关于无穷小微积分符号的工作开始。他本人完全意识到这一事实，并明确地将他的指数记号和行列式思想（{[}12 a{]}，第II卷，第240页；参见{[}12 bis{]}，第481-487页）以及他的“几何演算”草图（参见{[}12 bis{]}，第IX章）与他的“特征”联系起来。在他心中，其方案的核心部分应当是符号逻辑，或者用他的话来说，是一种“推理演算”；如果他未能成功创建这种演算，他至少做了三次尝试。在第一次尝试中，他想到将每个“原始”术语与一个素数相关联，由多个原始术语组成的每个术语则用相应素数的乘积来表示（*）；他试图将通常的三段论规则翻译到这个系统中，但遇到了由否定引起的相当大的复杂性（他自然试图用符号的变化来表示否定），并很快放弃了这种方法（{[}12 c{]}，第42-96页；参见{[}12 bis{]}，第326-344页）。在后来的尝试中，他试图赋予亚里士多德逻辑更代数的形式。有时他保留记号AB来表示两个概念的合取，有时他使用记号A + B（†）。他观察到（在乘法记号下）幂等律AA = A，指出命题“每个A都是B”可以用等式A = AB代替，并且从这个起点出发，亚里士多德的大部分规则可以通过纯代数计算获得（{[}12 c{]}，第229-237页和356-399页；参见{[}12 bis{]}，第345-364页）。他还有空概念（“非存在”）的想法，并认识到，例如，命题“每个A都是B”与“A.（非B）不成立”是等价的（同上）。此外，他注意到他的逻辑演算不仅适用于概念逻辑，也适用于命题逻辑（{[}12 c{]}，第377页）。因此，他已非常接近“布尔演算”。遗憾的是，他似乎未能完全摆脱经院哲学的影响；他不仅几乎完全将其演算的目标局限于用他的记号（**）转写三段论的规则，而且甚至不惜牺牲他最精妙的构想，以求完整恢复亚里士多德的规则，甚至包括那些与空集概念不相容的规则（††）。

莱布尼茨的著作大部分直到二十世纪初才得以出版，因此直接影响甚微。在整个十八世纪和十九世纪初，多位作者（德·塞格纳、兰伯特、普卢凯、霍兰德、德·卡斯蒂永、热尔戈纳）提出了与莱布尼茨类似的尝试纲要，但从未真正超越他所止步之处。他们的著作几乎没有产生什么影响，而且他们中的大多数人对前人的成果也一无所知（*）。在同样条件下写作的还有G. 布尔，他可被视为现代符号逻辑的真正创立者 {[}16{]}。他的核心思想是系统地采取“外延”观点，从而直接对集合进行运算，用 \(xy\) 表示两个集合的交，用 \(x + y\) 表示它们的并，当 \(x\) 与 \(y\) 不相交时。他还引入了“全域”，记为1（所有对象的集合），以及空集，记为0，并写作1 — \(x\) 来表示 \(x\) 的补集。正如莱布尼茨所做的那样，他将包含关系解释为关系 \(xy = x\) （由此他轻易地推导出古典三段论规则的合理性），而他对并和补的记号赋予了他的系统一种前所未有的灵活性（†）。此外，通过将每个命题与使其成立的所有“情形”的集合相关联，他将蕴涵关系解释为包含关系，他的集合演算由此为他提供了“命题演算”的规则。

在十九世纪下半叶，布尔系统成为一个活跃的逻辑学家学派工作的基础，他们在各个方面对其进行了改进和完善。例如，杰文斯（1864年）拓宽了并运算的意义 \(x + y\) ，将其推广到 \(x\) 与 \(y\) 为任意集合的情形。A. 德·摩根于1858年，C. S. 皮尔士于1867年证明了对偶关系

\[
\left(\mathbb {C} \mathrm{A}\right) \cap \left(\mathbb {C} \mathrm{B}\right) = \mathbb {C} (\mathrm{A} \cup \mathrm{B}), \quad \left(\mathbb {C} \mathrm{A}\right) \cup \left(\mathbb {C} \mathrm{B}\right) = \mathbb {C} (\mathrm{A} \cap \mathrm{B}) \quad (* *).
\]

1860年，德·摩根还开始了对关系的研究，定义了二元关系的逆和复合（即对应于图上的运算 \(\overline{\mathbf{G}}\) 并且 \(\mathbf{G}_1 \circ \mathbf{G}_2\) 的那些运算）（ \(\ddagger\) ）。所有这些工作都在施罗德卷帙浩繁的著作中得到了系统的阐述和发展 {[}23{]}。但颇为奇特的是，我们刚才提到的这些逻辑学家似乎对将其成果应用于数学几乎不感兴趣，相反，布尔，尤其是施罗德的主要目标显然是通过模仿古典代数的方法并复制其问题（往往相当人为）来发展“布尔”代数。这种态度的原因无疑在于，布尔演算仍然不适合转写大多数数学论证（*），因此只是对莱布尼茨宏伟梦想的一个非常不完整的回应。构建更适合数学的形式体系——引入变量和量词，这归功于弗雷格 {[}24{]} 和C. S. 皮尔士 {[}22 b{]}，是朝这个方向迈出的关键一步——这是逻辑学家和数学家们的工作，与他们的前辈不同，他们主要对数学基础的应用感兴趣。

弗雷格的计划 {[}24 b 和 c{]} 是将算术建立在用他命名为 Begriffsschrift（概念文字）的符号体系形式化的逻辑之上；我们稍后将回到他定义自然整数的方法。他的出版物以极端精确和对概念分析的细致入微为特征，这使他引入了许多此后被证明在现代逻辑中具有重大意义的区分。例如，他是第一个区分命题的陈述与命题为真的断定之间、属于关系与包含关系之间、对象x与以x为唯一元素的集合 \(\{x\}\) 等的人。这种形式化的逻辑不仅包含数学意义上的“变量”，还包含代表未定关系且可被量化的“命题变量”，后来（通过罗素和怀特海的工作）成为元数学的基本工具。遗憾的是，弗雷格采用的记号缺乏启发性，排版复杂得惊人，且远离当时数学的通行用法；由此造成的难以阅读大大削弱了他对同时代人的影响。

皮亚诺的目标既更宏大也更务实。他开始出版一部完全用形式化语言写成的《数学公式集》，其中不仅包含数理逻辑，还包含数学最重要分支的所有成果。他在一群热心的合作者（瓦伊拉蒂、皮耶里、帕多阿、瓦卡、

维万蒂、法诺、布拉利-福尔蒂）证明了他所采用的体系之卓越。紧随当时的数学惯例，并引入许多精心挑选的缩写符号，他的语言还成功做到了相当易读，这主要归功于一个巧妙的方法，即用分隔点来替代括号 {[}29 c{]} 。皮亚诺的许多记号如今已被大多数数学家普遍使用：例如， \(\in\) , \(\supset\) （但不同于现今用法，其含义为“包含于”或“蕴含”（*））， \(\cup\) , \(\cap\) ，A---B（差集 \(a—b\) ，其中 \(a\in A\) 并且 \(b\in B\) 此外，该《公式集》首次对函数的一般概念、直接像（ \(\dagger\) ）以及逆像作了仔细分析，并指出序列不过是定义在自然数集 N 上的函数。但在皮亚诺手中，量词受到了令人厌烦的限制（在他的体系中，原则上，只有形如 A \(\Longrightarrow\) B, A \(\Longleftrightarrow\) B，或A=B可量化）。此外，他某些门徒近乎狂热的热情使他极易成为笑柄，而庞加莱尤其经常提出不无偏颇的批评，这对皮亚诺学派是沉重打击，并阻碍了其学说在数学界的传播。

在弗雷格和皮亚诺那里，我们拥有了当今所使用的形式化语言的基本要素。这类语言中最著名的，毫无疑问，是罗素和怀特海在他们的巨著《数学原理》中所创造的那种，它巧妙地将弗雷格的精确性与皮亚诺的便捷性结合起来。当前使用的大多数形式化语言彼此之间的差异，仅在于次要的修改，其目的在于简化使用。其中最巧妙的一种是关系的“函数式”记号（例如， \(\in xy\) 以代替 \(x \in y\) ），这归功于卢卡西维茨，它完全省去了括号的需要。但最有趣的无疑是希尔伯特引入的符号 \(\tau\) ，这使得人们可以将量词 \(\exists\) 与 \(\forall\) 视为缩写符号，以避免皮亚诺和罗素的“全称”函数符号 \(\iota\) （该符号仅适用于函数关系），并在集合论中省去选择公理（{[}31{]}，第183页（ \(^{**}\) )).

